\documentclass[11pt]{article}

\usepackage[a4paper,margin=1in]{geometry}
\usepackage{amsmath,amssymb,amsthm,mathtools,bm}
\usepackage{enumitem}
\usepackage{booktabs}
\usepackage[hidelinks]{hyperref}
\usepackage[nameinlink,noabbrev]{cleveref}
\usepackage{microtype}

\newcommand{\E}{\mathbb{E}}
\newcommand{\R}{\mathbb{R}}
\newcommand{\cN}{\mathcal{N}}
\newcommand{\cL}{\mathcal{L}}
\newcommand{\W}{\mathcal{W}}
\newcommand{\dd}{\,\mathrm{d}}
\newcommand{\norm}[1]{\left\lVert #1 \right\rVert}
\newcommand{\data}{\mathrm{data}}

\theoremstyle{plain}
\newtheorem{theorem}{Theorem}
\newtheorem{proposition}{Proposition}
\newtheorem{lemma}{Lemma}

\theoremstyle{definition}
\newtheorem{definition}{Definition}
\newtheorem{remark}{Remark}

\title{Deep Theory Analysis for General Stochastic-Interpolant E-PVM}
\author{}
\date{\today}

\begin{document}

\maketitle

\begin{abstract}
This note analyzes the theoretical spine of endpoint-posterior velocity matching
(E-PVM) for generator-induced stochastic interpolants. The main message is that
the endpoint-posterior velocity identity is algebraically clean, but fixed-level
velocity matching alone is generally insufficient to guarantee endpoint
generation accuracy. For general stochastic interpolants, the conditional
velocity mixes the posterior mean of the base endpoint and the posterior mean of
the data endpoint. Thus a theory must either use a special bridge/base where the
base posterior can be eliminated, add a score or density closure equation, or
prove a new single-level velocity identifiability theorem under structural
assumptions.
\end{abstract}

\tableofcontents

\section{Main Takeaway}

For a general stochastic interpolant, the endpoint-posterior velocity identity
is clean, but fixed-level velocity matching alone is usually not enough to prove
endpoint generation accuracy. The reason is structural:
\begin{equation}
v_t(x)
=
\lambda_t x
+A_t m_t^0(x)
+B_t m_t^1(x).
\end{equation}
This is only one vector equation, while the posterior contains two unknown
endpoint posterior means:
\begin{equation}
m_t^0(x)=\E[X_0\mid I_t=x],
\qquad
m_t^1(x)=\E[X_1\mid I_t=x].
\end{equation}
Thus the theory must either:
\begin{enumerate}[label=(\roman*)]
\item choose a special bridge/base where \(m_t^0\) can be integrated out;
\item add an intermediate density or score control that gives a second equation;
\item prove a new single-level velocity identifiability theorem under strong
structural assumptions.
\end{enumerate}
This should be the central theoretical framing.

\section{General Identity}

Let the data interpolant be
\begin{equation}
I_t
=
a_tX_0+b_tX_1+\gamma_t\varepsilon,
\qquad
\varepsilon\sim\cN(0,I_D),
\end{equation}
and let the model interpolant be
\begin{equation}
I_{\theta,t}
=
a_tX_0+b_t h_\theta(U)+\gamma_t\varepsilon,
\qquad
U\sim\pi.
\end{equation}
Define
\begin{equation}
\lambda_t=\frac{\dot\gamma_t}{\gamma_t},
\qquad
A_t=\dot a_t-\lambda_t a_t,
\qquad
B_t=\dot b_t-\lambda_t b_t.
\end{equation}

\begin{proposition}[Endpoint-posterior velocity identity]
Assume \(\gamma_t>0\). The data conditional velocity satisfies
\begin{equation}
v_{\data,t}(x)
=
\lambda_t x
+A_t m_{\data,t}^0(x)
+B_t m_{\data,t}^1(x),
\end{equation}
where
\begin{equation}
m_{\data,t}^0(x)
=
\E[X_0\mid I_t=x],
\qquad
m_{\data,t}^1(x)
=
\E[X_1\mid I_t=x].
\end{equation}
The model conditional velocity satisfies
\begin{equation}
v_{\theta,t}(x)
=
\lambda_t x
+A_t m_{\theta,t}^0(x)
+B_t m_{\theta,t}^1(x),
\end{equation}
where
\begin{equation}
m_{\theta,t}^0(x)
=
\E_\theta[X_0\mid I_{\theta,t}=x],
\qquad
m_{\theta,t}^1(x)
=
\E_\theta[h_\theta(U)\mid I_{\theta,t}=x].
\end{equation}
\end{proposition}

\begin{proof}[Algebra]
On the data bridge,
\begin{equation}
\varepsilon
=
\frac{x-a_tX_0-b_tX_1}{\gamma_t}
\end{equation}
when conditioning on \(I_t=x\). Substituting this into
\begin{equation}
\dot I_t=\dot a_tX_0+\dot b_tX_1+\dot\gamma_t\varepsilon
\end{equation}
and taking the conditional expectation gives the claimed expression. The model
identity is identical with \(X_1\) replaced by \(h_\theta(U)\).
\end{proof}

\section{Why General SI Is Harder Than MAGT}

MAGT's Gaussian smoothing effectively has one posterior mean:
\begin{equation}
Y_t=\alpha_tX+\sigma_tZ,
\qquad
m_t(x)=\E[X\mid Y_t=x].
\end{equation}
The score satisfies a Tweedie identity:
\begin{equation}
s_t(x)
=
\nabla\log p_t(x)
=
-\frac{x-\alpha_t m_t(x)}{\sigma_t^2}.
\end{equation}
Therefore controlling the posterior mean is equivalent to controlling the score,
and the proof chain is:
\begin{equation}
\text{posterior mean error}
\Longleftrightarrow
\text{score/Fisher error}
\Longrightarrow
W_2(p_t,p_{\theta,t})
\Longrightarrow
W_2(P_X,P_{h_\theta(U)}).
\end{equation}
For a general stochastic interpolant, however,
\begin{equation}
v_t(x)
=
\lambda_t x+A_t m_t^0(x)+B_t m_t^1(x)
\end{equation}
does not directly give a Fisher divergence. It mixes the base posterior and the
endpoint posterior. This is the main theoretical obstruction.

\section{A Useful Second Equation: Score Closure}

The intermediate score is
\begin{equation}
s_t(x)
=
\nabla\log p_t(x)
=
-\frac{
x-a_t m_t^0(x)-b_t m_t^1(x)
}{
\gamma_t^2
}.
\end{equation}
Equivalently,
\begin{equation}
a_t m_t^0(x)+b_t m_t^1(x)
=
x+\gamma_t^2 s_t(x).
\end{equation}
Together with
\begin{equation}
A_t m_t^0(x)+B_t m_t^1(x)
=
v_t(x)-\lambda_t x,
\end{equation}
we obtain the two-equation system
\begin{equation}
\begin{pmatrix}
a_t & b_t\\
A_t & B_t
\end{pmatrix}
\begin{pmatrix}
m_t^0(x)\\
m_t^1(x)
\end{pmatrix}
=
\begin{pmatrix}
x+\gamma_t^2s_t(x)\\
v_t(x)-\lambda_t x
\end{pmatrix}.
\end{equation}
The determinant is
\begin{equation}
\Delta_t
=
a_tB_t-b_tA_t
=
a_t\dot b_t-b_t\dot a_t.
\end{equation}
Thus, if
\begin{equation}
\Delta_t\ne0,
\end{equation}
then the pair \((s_t,v_t)\) identifies both posterior means.

\begin{proposition}[Score--velocity closure]
If \(\Delta_t\neq0\), then
\begin{equation}
m_t^1(x)
=
\frac{
a_t\{v_t(x)-\lambda_t x\}
-A_t\{x+\gamma_t^2s_t(x)\}
}{
\Delta_t
}.
\end{equation}
Consequently,
\begin{equation}
\norm{m_{\theta,t}^1-m_{\data,t}^1}
\le
\frac{|a_t|}{|\Delta_t|}
\norm{v_{\theta,t}-v_{\data,t}}
+
\frac{|A_t|\gamma_t^2}{|\Delta_t|}
\norm{s_{\theta,t}-s_{\data,t}},
\end{equation}
for any common norm under which the displayed quantities are defined.
\end{proposition}

This is a useful bridge between general stochastic-interpolant velocity theory
and MAGT-style score theory:
\begin{equation}
\text{velocity mismatch}
+
\text{score/density mismatch}
\Longrightarrow
\text{endpoint posterior mismatch}.
\end{equation}

\section{Gaussian-Base Reduction}

Suppose
\begin{equation}
X_0\sim\cN(0,I_D),
\qquad
X_0\perp X_1,
\end{equation}
and
\begin{equation}
I_t=a_tX_0+b_tX_1+\gamma_t\varepsilon.
\end{equation}
Let
\begin{equation}
S_t=a_t^2+\gamma_t^2.
\end{equation}
Conditioned on \(X_1=y\),
\begin{equation}
I_t\mid X_1=y
\sim
\cN(b_ty,S_tI_D).
\end{equation}
Thus the bridge is a Gaussian corruption of the endpoint \(X_1\), with signal
coefficient \(b_t\) and noise variance \(S_t\). Moreover,
\begin{equation}
m_t^0(x)
=
\frac{a_t}{S_t}
\left(
x-b_tm_t^1(x)
\right).
\end{equation}
Substituting into the velocity identity yields
\begin{equation}
v_t(x)
=
D_t x+C_t m_t^1(x),
\end{equation}
where
\begin{equation}
D_t
=
\lambda_t+\frac{A_ta_t}{S_t},
\qquad
C_t
=
B_t-\frac{A_ta_tb_t}{S_t}.
\end{equation}
If
\begin{equation}
C_t\ne0,
\end{equation}
then
\begin{equation}
m_t^1(x)
=
\frac{v_t(x)-D_tx}{C_t}.
\end{equation}
In this special case, velocity matching is equivalent to endpoint posterior mean
matching. The score is also
\begin{equation}
s_t(x)
=
-\frac{x-b_tm_t^1(x)}{S_t}.
\end{equation}
Therefore,
\begin{equation}
v_{\theta,t}-v_{\data,t}
=
C_t
\left(
m_{\theta,t}^1-m_{\data,t}^1
\right)
=
-\frac{C_tS_t}{b_t}
\left(
s_{\theta,t}-s_{\data,t}
\right),
\end{equation}
assuming \(b_t\ne0\). This is the clean subcase where the theory can inherit
MAGT almost directly.

For the roadmap's linear noisy choice,
\begin{equation}
a_t=1-t,
\qquad
b_t=t,
\qquad
\gamma_t^2=2t(1-t),
\end{equation}
we have
\begin{equation}
S_t=1-t^2,
\end{equation}
and the velocity simplifies to
\begin{equation}
v_t(x)
=
\frac{m_t^1(x)-tx}{1-t^2}.
\end{equation}
This special case should be presented as the first tractable theorem, not as the
whole novelty.

\section{What a General SI Theory Must Prove}

\subsection{Theorem 1: Endpoint-Posterior Velocity Identity}

Prove
\begin{equation}
v_{\theta,t}(x)
=
\lambda_t x
+A_t m_{\theta,t}^0(x)
+B_t m_{\theta,t}^1(x).
\end{equation}
This is mostly algebra, but it establishes the conceptual object.

\subsection{Theorem 2: Population Velocity Calibration}

For any measurable \(v_{\theta,\tau}\),
\begin{align}
&\E
\left[
\norm{v_{\theta,\tau}(I_\tau)-\dot I_\tau}^2
\right] \nonumber\\
&\quad =
\E
\left[
\norm{v_{\data,\tau}(I_\tau)-\dot I_\tau}^2
\right]
+
\E_{p_\tau}
\left[
\norm{v_{\theta,\tau}-v_{\data,\tau}}^2
\right].
\end{align}
This says the velocity objective is calibrated to fixed-level conditional
velocity error.

\subsection{Theorem 3: Gaussian-Base Velocity--Score Equivalence}

Under Gaussian base and independent endpoint coupling,
\begin{equation}
\norm{v_{\theta,\tau}-v_{\data,\tau}}_{L^2(p_\tau)}^2
\asymp
\norm{s_{\theta,\tau}-s_{\data,\tau}}_{L^2(p_\tau)}^2,
\end{equation}
with explicit constants determined by \(a_\tau,b_\tau,\gamma_\tau\). This theorem
explains precisely why the Gaussian noisy case is close to MAGT.

\subsection{Theorem 4: General SI Closure With Score/Density Residual}

For a general base/interpolant, prove something like
\begin{align}
&\norm{m_{\theta,\tau}^1-m_{\data,\tau}^1}_{L^2(p_\tau)}
\nonumber\\
&\quad \le
C_\tau
\left(
\norm{v_{\theta,\tau}-v_{\data,\tau}}_{L^2(p_\tau)}
+
\gamma_\tau^2
\norm{s_{\theta,\tau}-s_{\data,\tau}}_{L^2(p_\tau)}
\right),
\end{align}
provided
\begin{equation}
\Delta_\tau=a_\tau\dot b_\tau-b_\tau\dot a_\tau
\ne0.
\end{equation}
This is not yet endpoint \(W_2\), but it is the missing algebraic closure.

\subsection{Theorem 5: Single-Level Velocity Pull-Back}

The ambitious target is
\begin{equation}
W_2(P_{X_1},P_{h_\theta(U)})
\le
C_\tau
\norm{v_{\theta,\tau}-v_{\data,\tau}}_{L^2(p_\tau)}
+
R_\tau.
\end{equation}
The residual \(R_\tau\) should be made explicit. Likely candidates are
\begin{equation}
R_\tau
=
C_\tau'
W_2(p_{\theta,\tau},p_\tau),
\end{equation}
or
\begin{equation}
R_\tau
=
C_\tau'
\norm{s_{\theta,\tau}-s_{\data,\tau}}_{L^2(p_\tau)}.
\end{equation}
The pure velocity-only version is the hardest and may be false without extra
assumptions.

\subsection{Theorem 6: Finite Estimator Perturbation}

For anchors or any estimator,
\begin{equation}
\E_{p_\tau}
\left[
\norm{\widehat v_{\theta,K,\tau}-v_{\theta,\tau}}^2
\right]
\le
\varepsilon_K^2.
\end{equation}
Then the population theorem receives an additive perturbation:
\begin{equation}
W_2(P_{X_1},P_{h_\theta(U)})
\le
C_\tau
\left(
\widehat{\mathcal E}_{\mathrm{vel}}^{1/2}
+
\varepsilon_K
\right)
+
R_\tau.
\end{equation}
This keeps anchors out of the main theorem until a concrete estimator is
selected.

\section{Recommended Proof Strategy}

\subsection{Stage A: Safe Theorem}

Prove the Gaussian-base reduction:
\begin{equation}
X_0\sim\cN(0,I_D)
\quad
\Longrightarrow
\quad
v_t(x)=D_tx+C_tm_t^1(x).
\end{equation}
Then show velocity matching is equivalent to score/denoiser matching and
inherits MAGT-style pull-back. This theorem is low risk but MAGT-adjacent.

\subsection{Stage B: General Identity and Closure}

Prove the two-equation closure:
\begin{equation}
(s_t,v_t)
\Longleftrightarrow
(m_t^0,m_t^1)
\quad
\text{if}
\quad
\Delta_t\ne0.
\end{equation}
This is genuinely general stochastic-interpolant theory and clarifies why
velocity alone is insufficient.

\subsection{Stage C: Partial Pull-Back}

Prove a theorem with an explicit residual:
\begin{equation}
W_2(P_{X_1},P_{h_\theta(U)})
\le
C_\tau
\norm{v_{\theta,\tau}-v_{\data,\tau}}_{L^2(p_\tau)}
+
C_\tau'
W_2(p_{\theta,\tau},p_\tau).
\end{equation}
This is easier and honest: fixed-level velocity needs intermediate distribution
alignment unless the Gaussian-base reduction applies.

\subsection{Stage D: Pure Velocity Pull-Back, If Possible}

Try to remove the residual using a growth lemma:
\begin{equation}
\norm{v_{\theta,t}-v_{\data,t}}_{L^2(p_t)}
\le
G(t,\tau)
\norm{v_{\theta,\tau}-v_{\data,\tau}}_{L^2(p_\tau)}.
\end{equation}
This likely requires strong assumptions:
\begin{equation}
\norm{\nabla v_t}_{\infty}\le L_t,
\qquad
\norm{\nabla^2\log p_t}_{\infty}\le H_t,
\end{equation}
plus nondegenerate density and manifold tube conditions.

\section{Bridge-Design Insight}

The coefficients
\begin{equation}
A_t=\dot a_t-\lambda_ta_t,
\qquad
B_t=\dot b_t-\lambda_tb_t
\end{equation}
determine how much the velocity depends on base posterior versus endpoint
posterior. If one can choose a schedule with small \(|A_\tau|\), then
\begin{equation}
v_{\theta,\tau}(x)
\approx
\lambda_\tau x+B_\tau m_{\theta,\tau}^1(x),
\end{equation}
which makes endpoint recovery easier. The condition
\begin{equation}
A_\tau=0
\end{equation}
means
\begin{equation}
\frac{\dot a_\tau}{a_\tau}
=
\frac{\dot\gamma_\tau}{\gamma_\tau}.
\end{equation}
This may be impossible globally with endpoint boundary conditions, but it can
guide a single-level bridge design. This is a potential theoretical design knob
distinct from anchors.

\section{Suggested Revised Contribution Statement}

The paper should claim:
\begin{enumerate}
\item We derive endpoint-posterior velocity identities for generator-induced
noisy stochastic interpolants.
\item We show that the Gaussian-base noisy bridge reduces to a MAGT-like
endpoint posterior mean and admits velocity--score equivalence.
\item For general stochastic interpolants, we prove that velocity alone mixes
base and data endpoint posteriors; a score/density closure or special bridge
design is required.
\item We formulate and partially prove a single-level velocity pull-back theorem
with explicit residual terms.
\item We treat anchors as a generic plug-in estimator satisfying an \(L^2\)
velocity approximation condition.
\end{enumerate}

This is a cleaner and more defensible theory story than making anchor design the
primary novelty.

\end{document}
